\documentclass[11pt]{article}
\usepackage[margin=1in]{geometry}
\usepackage{amsmath,amssymb,amsthm,mathtools}
\usepackage[hidelinks]{hyperref}
\newtheorem{theorem}{Theorem}[section]
\newtheorem{lemma}[theorem]{Lemma}
\newtheorem{proposition}[theorem]{Proposition}
\newtheorem{remark}[theorem]{Remark}
\newcommand{\E}{\mathbb E}
\newcommand{\dd}{\,d}
\newcommand{\one}{\mathbf1}
\title{Residual closure and small-constant Q2 transport}
\author{Analytical proof increment for review}
\date{Version 1; October 5, 2026}
\begin{document}\maketitle
\begin{abstract}
We extend BULK's conditional ratio bootstrap to a small constant entry
ratio and prove the target-only threshold dichotomy with its necessary
endpoint regularization. We also prove a narrower residual estimate
conditional on an independently controlled clock. Stage 1 is not
closed: the scalar fixed point proposed in S0 drops a regenerated-energy
term, and the attempted Volterra repair cannot yet remove the clock
premise without a circular argument. We display both failures and stop;
Stage 2 is not started. No original-flow entry, capped gain, or
strong-learning estimate is asserted here.
\end{abstract}

\section{Conventions and scope}
All labels in this increment have prefix \texttt{an03rcetv1:}.
Use $P_1=N(e_1,q^2I)$, $P_2=N(\rho e_2,q^2I)$,
$P=(P_1+P_2)/2$, $\rho\in[13/20,7/10]$, $\lambda=\rho^2$,
$Q=\lambda+q^2$, and fixed initial label measure $\dd\lambda_0$.
The original equations are
\[
 U'=R_\theta^TW,\quad W'=R_\theta U,\quad
 R_\theta=\E_P[(X-f(X))X^T\one_{\{a_\theta^TX>0\}}],\quad |U|=|W|.
\]
For a fixed cohort, use BULK's $H,A,E=H+A,J$, with
$z=(U_2+W_2)/2$, $a=(U_2-W_2)/2$ and
$j=((U_1^2+W_1^2)/2)^{1/2}$. Thus $H=\int z^2$,
$A=\int a^2$, and $J=\int j^2$.
The response is $c=\E_{P_2}[X_2f_2]/Q$, and $I'=Q(1-c)$.
Constants may depend on fixed source-contract constants, but not on
$q$ or the labels. Original-flow conclusions below are conditional
implications removing residual hypotheses; the target-only result has
its separate exact-comparison scope.

\section{The small-constant bootstrap}
\begin{theorem}[Conditional original-flow transport]
\label{an03rcetv1:eta0}
Retain every hypothesis of BULK Theorem 4.1 except its entry power:
on $[t_0,t_b]$, $0\le c\le c_b<1$, $t_b-t_0\le L_T\log(1/q)$,
the diagonal residual norms are bounded by $M$, the sum of cross
residual norms is bounded by $Mq$, and
$\int r\le B$, where $r=\operatorname*{ess\,sup}_G|R_{11}|$.
At entry suppose $z>0$, $|a|/z\le1/4$ and $qj/z\le C_0$.
There is $C_{\rm adm}>0$, depending only on the stated constants,
such that for $0<C_0<C_{\rm adm}$ and sufficiently small $q$
the following hold. With $b_*=Q(1-c_b)/2$ and $\tau=t-t_0$,
\begin{align}
 x:=qj/z&\le C(C_0e^{-b_*\tau/2}+q^2),
 \qquad |a|/z\le1/3, \label{an03rcetv1:ratio}\\
 \frac{a_{\theta,2}}q&\ge c\min\{q^{-1},C_0^{-1}e^{b_*\tau/2}\},
 \label{an03rcetv1:margin}\\
 \int_s^t|\varepsilon_\alpha(v)|\dd v,
 \ \int_s^t|\varepsilon_G(v)|\dd v
 &\le C\{C_0e^{-b_*(s-t_0)/2}+q^2(t-s)\},
 \label{an03rcetv1:defect}
\end{align}
where $\varepsilon_\alpha=(\log z_\alpha^2)'-I'$ and
$\varepsilon_G=(\log H_G)'-I'$. In particular
$D_{\rm osc}\le C(C_0+q^2\log(1/q))$.
If $G_1,G_2$ are positive-amplitude fixed subsets of the guarded
cohort and $D=C(C_0+q^2\log(1/q))$, then
\begin{equation}\label{an03rcetv1:transport}
 e^{-2D}\le
 \frac{H_{G_2}(t)/H_{G_1}(t)}{H_{G_2}(t_0)/H_{G_1}(t_0)}
 \le e^{2D},\qquad H_{G_i}\le E_{G_i}\le(10/9)H_{G_i}.
\end{equation}
\end{theorem}
\begin{proof}
We use the exact differential inequalities and Gaussian estimate
already proved in BULK, rather than reprove their coordinate algebra.
For a fixed $L$ supplied by its residual-norm bounds, $y=|a|/z$ and
$\Delta=|R_{22}-Q(1-c)/2|$ satisfy, while $z>0$,
\begin{align*}
 D^+x&\le(r-b+\Delta+Lx)x+Lq^2\sqrt{1+y^2},\\
 D^+y&\le(-2b+2\Delta+Lx)y+Lx,
 \qquad z'/z\ge b-\Delta-Lx,\\
 \Delta&\le Cq^2+C\exp[-\rho^2a_{\theta,2}^2/(8q^2)],
 \qquad b=Q(1-c)/2\ge b_*.
\end{align*}
Take the fixed guard $X=4e^BC_0$, $x<X$, $y<1/2$, $z>0$.
Under it, $U_2\ge z/2$ and
$|q\cot\theta|\le2\sqrt2x$, hence
$a_{\theta,2}/q\ge(q^2+8X^2)^{-1/2}$.
Choose $C_0$ first so that $LX\le b_*/4$ and the resulting
Gaussian tail is at most $b_*/16$, uniformly for small $q$.
Then decrease $q$ so that $\Delta\le b_*/8$. The inequalities give
\[
 D^+x\le(r-b_*/2)x+2Lq^2,
 \quad x(t)\le e^BC_0e^{-b_*\tau/2}+4Le^Bq^2/b_*<X/2.
\]
Also $D^+y\le-3b_*y/2+Lx$ gives
$y\le1/4+LX/(3b_*)<1/3$ after decreasing $C_0$.
Finally $z'/z\ge b_*/2$, excluding a zero of $z$.
All first-exit guards improve strictly.

The improved $x$ estimate, not the initial constant guard, now gives
\eqref{an03rcetv1:margin}. More precisely, with fixed $A,c>0$,
\[
 \Delta\le Cq^2+C\exp\!\left[-\frac{c}{q^2+A^2C_0^2e^{-b_*\tau}}\right]
 \le Cq^2+C C_0e^{-b_*\tau/2}.
\]
For the second inequality split the exponential using
$(u+v)^{-1}\ge\frac12\min(u^{-1},v^{-1})$ and
$e^{-c/w^2}\le C_c w$ for $w>0$; the $q$-only exponential
is at most $Cq^2$. These are analytic inequalities, uniform in time.
Thus the initially constant gate error has an integrable tail.
The exact $z$ equation yields
$|\varepsilon_\alpha|\le2\Delta+2Lx$, proving
\eqref{an03rcetv1:defect}. The aggregate defect is its current
$z_\alpha^2$-weighted average. Integrating the difference of two
aggregate clock identities proves \eqref{an03rcetv1:transport}.
Only $z>0$ is divided by; derivatives of $j$ at zero are norm Dini
derivatives. Finite coordinate energy supplies domination for the
fixed-cohort integrals.
\end{proof}

\begin{remark}\label{an03rcetv1:clock}
For fixed $C_0$, this gives bounded transport distortion, not an
$o(1)$ clock defect. Retain BULK's $\eta=1/8$ cohort
$\mathcal C_q=\{c_0\ge q^{1/4}\}$ as the phase clock.
No $o(1)$ phase conclusion follows from \eqref{an03rcetv1:defect} alone.
\end{remark}

\section{The target-only threshold and its endpoint correction}
Set $T=10\log(1/q)$, $\omega=(1-\lambda)/2$,
$r_\lambda=(1-\lambda)/\lambda$, and
$\gamma=1-5\lambda+2\lambda/(1-\lambda)$.
For an initial Q2 label $\alpha$, put $c_0=-\cos\alpha$,
$s_0=\sin\alpha$. Hats refer exclusively to the aligned target-only
comparison.
\begin{proposition}[Exact comparison dichotomy]
\label{an03rcetv1:dichotomy}
For labels not yet in a fixed sufficiently large strong chart at $T$,
\begin{equation}\label{an03rcetv1:regularized}
 q\cot\widehat\theta(T,\alpha)
 \asymp\left[q^\gamma\frac{s_0+q}{c_0+q}\right]^{r_\lambda}.
\end{equation}
In the threshold region $c_0\asymp q^\gamma$, and more generally
when $c_0\gg q$ and $s_0$ is bounded below, this becomes
\begin{equation}\label{an03rcetv1:thresholdpower}
 q\cot\widehat\theta(T,\alpha)
 \asymp(q^\gamma/c_0)^{r_\lambda}.
\end{equation}
Given any fixed small $C_0>0$, choose a fixed $K$ sufficiently large.
Every label with $c_0\ge Kq^\gamma$ has, at $T$,
$q\widehat j/\widehat z\le C_0$ and $\widehat a=0$.
Every remaining label has both aligned scaled coordinates in a
fixed strong chart, with displacement from $a_q$ bounded by a
constant $Z=Z(C_0)$ independent of $q$.
\end{proposition}
\begin{proof}
Q2E's uniform crossing formula gives
\[
 T_\times=\frac2\lambda\log\frac{c_0+q}{q(s_0+q)}+O(1),
 \quad \sup T_\times=\frac4\lambda\log(1/q)+O(1).
\]
Hence $v=T-T_\times\to\infty$ uniformly, since
$4/\lambda\le1600/169<10$. After a fixed positive layer time,
Q2E's created first-coordinate seed and its upper bound imply
\[
 \widehat U_1(T_\times+v)\asymp
 q\widehat U_2(T_\times)e^{v/2}.
\]
The lower bound follows by retaining the positive seed and the
exact inequality $\widehat U_1'\ge\widehat U_1/2$.
Before arrival at $qL$, the extra second-coordinate logarithmic
rate above $\lambda/2$ has bounded integral. Indeed in the part
$\theta\ge\pi/4$ it is exponentially small over an $O(\log(1/q))$
time, and on $qL\le\theta\le\pi/4$ PROFILE's reciprocal-speed
estimate bounds it by
$C\int_{L/2}^\infty K(-\rho z)z^{-2}\dd z$.
The $d_q$ contribution is $O(q^2\log(1/q))$.
Consequently $\widehat U_2\asymp
\widehat U_2(T_\times)e^{\lambda v/2}$ there. Taking the ratio
and inserting $T_\times$ proves \eqref{an03rcetv1:regularized}, since
$2-10\omega+2\omega/\lambda=\gamma r_\lambda$.

For the large-$c_0$ group the one-sided estimate
$q\cot\widehat\theta(T)\le Cq^2e^{\omega(T-T_\times)}$
holds even without a pre-arrival assumption (BULK's proof of its
bulk-seed theorem). It is at most $CK^{-r_\lambda}$.
Choose $K$ so this is at most $C_0$. Here $\widehat j=\widehat U_1$
and $\widehat z=\widehat U_2$, so no alignment factor is needed.
For the other group, $c_0\le Kq^\gamma$ implies $s_0\asymp1$,
and the crossing formula gives
\[
 T-T_\times\ge(2/\omega)\log(1/q)-C_K.
\]
PROFILE's axis-start arrival time differs from the first term by
$O(1)$. At most a fixed time remains until arrival. On the reverse
segment above $qL$, its outer decomposition gives
$|\theta'|\le C\theta$ on $[qL,\pi/4]$. Starting at $qL$ and
integrating backwards for this fixed time keeps $\theta/q$ bounded;
for small $q$ it cannot reach $\pi/4$. Already-arrived labels stay
in the invariant chart. This proves the fixed cap $Z$.
\end{proof}

\begin{remark}\label{an03rcetv1:endpoint}
The unregularized two-sided formula is not uniform over all Q2.
At $\alpha=\pi$, \eqref{an03rcetv1:regularized} contains the extra
factor $q^{r_\lambda}$; this label is still outside the strong chart
at $T$. The threshold dichotomy is unaffected because near its
cutoff $s_0\asymp1$ and $c_0\gg q$. The proposition supplies
comparison entry geometry only. Applying Theorem
\ref{an03rcetv1:eta0} or BULK Proposition 5.2 to original labels still
requires original joint-entry transfer and, for Proposition 5.2,
its pre-exit distance and capped-gain contracts. A fixed displacement
cap does not prove those contracts.
\end{remark}

\section{Residual closure from the stated output budgets}
Let fixed disjoint groups $K,G,B$ partition the population on the
passage interval: a strong core, the explicit-output cohorts (in
particular the whole Q2 ancestry), and the outside population.
Use the following existing S0/BULK input budgets:
\begin{align}
 |\theta_\alpha|+|\psi_\alpha|&\le Rq\quad(\alpha\in K),
 \quad \mu_K\le M_K,\quad
 \sup_t\mu_B(t)/q=o(1), \label{an03rcetv1:coreinputs}\\
 E_G(t)&\le C_E(q^5+H(t)),\quad J_G(t_0)\le C_Jq^3,
 \label{an03rcetv1:energyinputs}\\
 H(t_b)&\le H_b,\quad
 H(s)\le e^D H(t)e^{-b_0(t-s)}\quad(s\le t),
 \quad b_0=Q(1-c_b),\quad D\le1, \label{an03rcetv1:clockinput}\\
 \kappa_1(t)&=\int_K W_1U_1\dd\lambda_0,\quad
 \int_{t_0}^{t_b}|1-\kappa_1|\dd t\le B_s.
 \label{an03rcetv1:learninginput}
\end{align}
Here $H$ is the $\eta=1/8$ clock. The backward estimate
\eqref{an03rcetv1:clockinput} is precisely the consequence of
BULK Theorem 4.1 used in its Theorem 6.2. Here it must be an
\emph{independently supplied premise}: deriving it from the residual
conclusion below would be circular. This retained clock premise is
why the following theorem does not complete the requested S0 closure.
The energy budget is an existing
conditional input. It follows from primary relative-amplitude
transport with an entry-energy comparison to the clock, together
with BULK's secondary-energy implication if its contracts hold.
No derivation of these entry or gain contracts is included here.

\begin{lemma}[Instantaneous output bounds and the strong remainder]
\label{an03rcetv1:outputs}
Under \eqref{an03rcetv1:coreinputs}, write $f_1=\kappa_1X_1+g_1$
on $P_1$ and put $Y=\sqrt{J_G}/q$. Uniformly over receivers,
\begin{align}
 \frac{\|e_2\|_{P_1}+\|e_1\|_{P_2}}q
 &\le C\{1+E_G+q^2Y^2+\sqrt{E_G}Y+\mu_B/q\},
 \label{an03rcetv1:crossnorm}\\
 \|g_1\|_{P_1}&\le C\{q^2+q^2Y^2+q^2Y\sqrt{E_G}+\mu_B\}.
 \label{an03rcetv1:g1}
\end{align}
Here $\|\cdot\|_{P_i}$ denotes the $L^2(P_i)$ norm.
\end{lemma}
\begin{proof}
The target cross-coordinate norms are $q$. For a core donor the
second output component is $O(q\sqrt m)$; on $P_2$ its input
feature has norm $O(q\sqrt m)$. Minkowski therefore bounds each
core cross output by $Cq\mu_K$.
On $P_1$, replace the core ReLU by its linear feature. The coefficient
of $X_2$ is $\int_K W_1U_2=O(q\mu_K)$ and $\|X_2\|_{P_1}=q$.
The replacement error is exponentially small: the core input
projection has mean bounded below and variance $q^2$ after dividing
by $\sqrt m$, and its negative-part second moment is bounded by
Gaussian fourth moments times the square root of its tail probability.
It is at most $Cq^2\mu_K$ for small $q$.

BULK's exact cross-output inequality is
$\mathfrak F_G\le8(E_G+J_G+\sqrt{E_GJ_G}/q)$.
Also Minkowski and Cauchy--Schwarz give directly
\[
 \|(f_G)_1\|_{P_1}
 \le\sqrt{1+q^2}\int_G|W_1U_1|
       +q\int_G|W_1U_2|
 \le C(J_G+q\sqrt{J_GE_G}).
\]
Finally $|f_B(X)|\le\mu_B|X|$ bounds every outside norm by
$C\mu_B$. Summing proves both claims. These estimates involve no
moving label sets or cancellation assumptions.
\end{proof}

\begin{theorem}[Narrow residual estimate with an independent clock bound]
\label{an03rcetv1:closure}
Under the budgets above and the logarithmic duration bound,
diagonal residual norms are uniformly bounded, cross residual norms
are $O(q)$, and
\begin{equation}\label{an03rcetv1:r11}
 \int_{t_0}^{t_b}\sup_\theta|R_{11}(t,\theta)|\dd t
 \le C\int_{t_0}^{t_b}|1-\kappa_1(t)|\dd t+o(1).
\end{equation}
There is no smallness requirement on the fixed cap $H_b$.
The independent clock hypothesis \eqref{an03rcetv1:clockinput}
is essential to the scope of this statement; its removal is not proved.
\end{theorem}
\begin{proof}
Loss monotonicity gives $\E_{P_i}|e|^2\le4\mathcal L(0)$.
The initialized loss is uniformly bounded, so the diagonal estimates
are immediate. BULK's residual-entry inequalities and
\eqref{an03rcetv1:crossnorm} give
\[
 \ell/q\le C(1+E+q^2Y^2+\sqrt E Y+\mu_B/q),
 \quad \ell=\sup_\theta\max(|R_{12}|,|R_{21}|).
\]
Its $R_{11}$ inequality, with $e_1=(1-\kappa_1)X_1-g_1$ on $P_1$,
also gives
\begin{equation}\label{an03rcetv1:rpoint}
 r\le C|1-\kappa_1|+
 C\{q^2(1+E+Y^2+\sqrt E Y)+\mu_B\},\qquad
 r=\sup_\theta|R_{11}|.
\end{equation}
From the backward clock bound and the energy budget,
\[
 \sup E\le C_*,\quad \int E\le K_E,\quad
 \int\sqrt E\le K_{1/2},
\]
with fixed constants, since
$\int H\le eH_b/b_0$, $\int\sqrt H\le2\sqrt{eH_b}/b_0$,
and $q^{5/2}(t_b-t_0)=o(1)$.

Here is the fixed point with its constants retained. Stop at $Y<K$
and $\int r<B$, where $B=CB_s+1$ with the coefficient from
\eqref{an03rcetv1:rpoint}. BULK's exact transverse inequality gives
\begin{equation}\label{an03rcetv1:volterra}
 Y(t)\le A_*+C e^B\int_{t_0}^t E(s)Y(s)\dd s
       +C e^Bq^2\int_{t_0}^t\sqrt{E(s)}Y(s)^2\dd s.
\end{equation}
One can take a fixed $A_*>1$ dominating
$e^B[\sqrt{C_Jq}+C\int\sqrt E(1+E+\mu_B/q)]$ for all small $q$.
Choose $K=4A_*\exp(Ce^BK_E)$. For sufficiently small $q$ the last
term under the guard is at most $A_*$, so Gronwall gives $Y\le K/2$.
Equation \eqref{an03rcetv1:rpoint} then gives $\int r\le CB_s+o(1)<B$.
Indeed $q^2\log(1/q)=o(1)$ and
$\sup\mu_B\log(1/q)=o(q\log(1/q))=o(1)$.
Both guards close, proving bounded $Y$, the cross estimate, and
\eqref{an03rcetv1:r11}.

This proof uses the global backward bound as an input even when
the $Y,r$ guards are stopped. Consequently it contains no step
claiming to close that clock input simultaneously.
\end{proof}

\begin{remark}[Why the direct scalar shortcut is insufficient]
\label{an03rcetv1:scalar}
Keeping the trial bound $\ell\le Lq$ in the transverse estimate yields
\[
 \sqrt J\lesssim q^{3/2}+Lq\sqrt H,
 \qquad \mathfrak F_G\lesssim
 q^3+H+\sqrt{qH}+LH+L^2q^2H.
\]
Thus a scalar supremum fixed point has the term $CLH_b$, which
small $q$ does not absorb. The expression
$H_b+L\sqrt{qH_b}+q^3$ in the proposed S0 shortcut omits this term.
It is not a valid consequence of BULK's regenerated-energy estimate
with constants uniform in the trial $L$. In
\eqref{an03rcetv1:volterra}, that same feedback is retained as the
integrable coefficient $CE(t)$. The resulting exponential depends
on $H_b$ but is finite for every fixed $H_b$. Only the remaining
quadratic term needs small $q$. This repairs the feedback estimate
conditional on clock control; it does not supply that control.
\end{remark}

\section{Obstruction: two unsuccessful closure routes}
\label{an03rcetv1:obstruction}
\paragraph{First route: a scalar constant fixed point.}
The proposed absorption would need
\[
 L_{\rm new}\le A+C H_b L+o_q(1)L<L.
\]
The last inequality does not follow by decreasing $q$: it would
require additional control such as $CH_b<1$. No such smallness
is supplied. Remark \ref{an03rcetv1:scalar} identifies the retained
term. Its scale is real at the level of the available energy data.
For example, a balanced aligned donor with
$U=W=(qL\sqrt E,\sqrt E)$ has $J=q^2L^2E$ and
\[
 \|(f_G)_2\|_{P_1}/q
 \ge\E_{P_1}(f_G)_2/q\ge LE,
\]
by convexity of the positive part. Thus the cross-output estimate
cannot uniformly replace this scale by $E+L\sqrt{qE}$ for arbitrary
trial $L$. This is an instantaneous-state example, not a population
trajectory or an initialized counterexample.

\paragraph{Second route: Volterra feedback with a stopped clock.}
Equation \eqref{an03rcetv1:volterra} produces a uniform constant only
if $\int E$, hence its supplied upper budget in terms of $\int H$,
is uniformly bounded. On a stopped interval ending at $t_*<t_b$,
BULK gives only
\[
 \int_{t_0}^{t_*}H(s)\dd s\le e H(t_*)/b_0.
\]
The required next inequality is
\begin{equation}\label{an03rcetv1:failingendpoint}
 H(t_*)\le e H_b\quad\text{uniformly at every first-exit endpoint}.
\end{equation}
The supplied cap is $H(t_b)\le H_b$. Inferring
\eqref{an03rcetv1:failingendpoint} from it requires the clock defect
bound on $[t_*,t_b]$, the as-yet-uncontrolled future interval.
Imposing a guard $H\le eH_b$ does not improve that guard: positivity
of $H'$ on the controlled past does not prevent a later fall to the
capped endpoint after other guards have failed.

The scalar endpoint-to-integral implication is false without clock
control. Let $\mathcal T=L_T\log(1/q)$, fix $k>0$, and define
\[
 H(t)=
 \begin{cases}
 H_b\exp[(2k/L_T)(t-\mathcal T/2)],&0\le t\le\mathcal T/2,\\
 H_b,&\mathcal T/2\le t\le\mathcal T.
 \end{cases}
\]
Then $H(0)=H_bq^k$, $H(\mathcal T)=H_b$, but
$\int H\ge(H_bL_T/2)\log(1/q)$.
With constant permitted $c$, the exact scalar clock identity holds
almost everywhere with $\varepsilon=(\log H)'-Q(1-c)$.
The plateau has a negative defect of logarithmic size.
This is a counterexample to the scalar inference, not a claimed
trajectory of the original flow. Loss monotonicity supplies diagonal
residual bounds, not a bound on this particular symmetric amplitude:
CLOCK's exact ledger $H=c-c_O-\tau+A$ retains signed complementary
response, antisymmetric energy, and a negative-gate term.

An independent uniform amplitude/response-ledger bound, an independent
all-subinterval clock estimate, or an appropriate signed transverse
estimate could repair the missing step. None is silently supplied
here. These two routes do not establish S0 from the requested inputs,
so the stop rule applies and Stage 2 is not begun. This obstruction
does not assert that the desired original-flow estimate is false.

\section{Precisely retained inputs}
The narrow theorem uses the following already open budgets:
(i) the original fixed-core chart and mass bounds and outside
$\mu_B=o(q)$; (ii) original Cartesian entry for the clock and
primary labels, their entry-amplitude comparison, and
$J_G(t_0)=O(q^3)$; (iii) the whole explicit-output energy bound
$E_G\le C(q^5+H)$, including all primary labels and secondary labels
(the latter still requires the transferred finite-cutoff profile,
pre-exit distance envelope, and capped gain of BULK Proposition 5.2);
(iv) the response guard, logarithmic duration and amplitude cap;
and (v) $\int|1-\kappa_1|\le B_s$.
It additionally retains the independent backward clock bound
\eqref{an03rcetv1:clockinput}. Removing that premise is unresolved.
None of (i)--(v) is claimed from initialization here. In particular
the comparison dichotomy does not silently transfer entry geometry.
No small terminal amplitude, raw $J=O(q^3)$ passage bound, or
vanishing defect for the $\eta=0$ transport cohort is claimed.

\begin{thebibliography}{9}
\bibitem{BULK} BULK, \textit{AN03\_bulk\_clock\_and\_secondary\_Q2\_energy\_v1}.
Source prefix \texttt{inc:AN03:bulkclock:v1:}; labels
\texttt{bulkclockclosure}, \texttt{R22tail}, \texttt{bulkweights},
\texttt{gainenergy}, \texttt{secondary}, \texttt{transverse},
\texttt{forcing}, \texttt{residualinputs}.
\bibitem{Q2E} Q2E, \textit{AN02\_Q2\_boundary\_layers\_and\_cohort\_energy\_v1}.
Prefix \texttt{inc:AN02:q2completion:v1:}; labels
\texttt{crossing}, \texttt{axisU1upper}, \texttt{profile}.
\bibitem{PROFILE} PROFILE, \textit{AN02\_strong\_entry\_profile\_and\_tail\_budget\_v1}.
Prefix \texttt{inc:AN02:profile:v1:}; labels \texttt{root},
\texttt{arrival}, \texttt{outerdecomposition}, \texttt{cost}.
\bibitem{F} \textit{RELU\_POPULATION\_FOUNDATIONS\_v2}, P5:
loss dissipation, balance, and characteristic regularity.
\bibitem{CLOCK} CLOCK,
\textit{AN03\_exact\_amplitude\_clock\_and\_c\_passage\_v1},
prefix \texttt{inc:AN03:clock:v1:}; labels
\texttt{endpoint}, \texttt{Hcap}, \texttt{crossenergy}.
\end{thebibliography}
\end{document}
